\section{Algebra of the four-weight construction}\label{app:candidate}
\label{alg:section}

Throughout, $\langle u,v\rangle=u^*v$, and a collection of $N$ vectors is
written as a matrix with $N$ columns. Put $\alpha=(1+\sqrt3)/4$.
A profile is admissible if $0<p_a<1/3$ for $1\leq a\leq4$ and
$\sum_a p_a=1$. Unsubscripted matrix norms are operator norms.
We distinguish the kernel projector $\Pi$, the mass matrix
$D_{\mathrm{mass}}$, the Bloch Gram difference $\Delta_{\mathrm B}$,
and the scalar candidate point bound $C_{\mathrm{pt}}$.
The local Hessian error bound in Section~\ref{app:local} is denoted
$E_{\mathrm{Hess}}$.

\begin{lemma}[The scalar parameter]\label{alg:root}
For every admissible profile there is exactly one
$\lambda_p\in(2/3,\alpha]$ satisfying
\begin{equation}\label{alg:root-equation}
 2\lambda_p=\sum_a\frac{p_a(1-3p_a)}{\lambda_p-2p_a}.
\end{equation}
Moreover $\lambda_p<1-\max_a p_a$, and $\lambda_p=\alpha$ if and only if
all four weights are $1/4$.
\end{lemma}
\begin{proof}
For $\lambda\geq2/3$ all denominators below are positive. Define
\begin{align}
 \Phi_p(\lambda)
 &=\lambda(3\lambda-2)\sum_a\frac1{\lambda-2p_a}
                  -2(2\lambda-1)\label{alg:Phi}\\
 &=4\left[2\lambda-\sum_a\frac{p_a(1-3p_a)}{\lambda-2p_a}\right].
 \nonumber
\end{align}
Polynomial division and $\sum p_a=1$ give the second identity. Thus
\[
 \Phi_p(2/3)=-2/3,\qquad
 \Phi_p'(\lambda)=4\left[2+
 \sum_a\frac{p_a(1-3p_a)}{(\lambda-2p_a)^2}\right]>8.
\]
Strict convexity gives $\sum_a(\alpha-2p_a)^{-1}\geq4/(\alpha-1/2)$.
Since $\alpha(3\alpha-2)>0$, substitution yields
$\Phi_p(\alpha)\geq0$, with equality precisely at the uniform profile.
Continuity and strict monotonicity prove existence and uniqueness.
For the extra bound put $p=\max p_a$ and $L=1-p>2/3$. The other three
weights sum to $L$, so their denominators $L-2p_a$ sum to $L$. Therefore
\[
 \sum_a\frac1{L-2p_a}\geq\frac1{3L-2}+\frac9L,\qquad
 \Phi_p(L)\geq L(3L-2)\left[\frac1{3L-2}+\frac9L\right]
                 -2(2L-1)=8(3L-2)>0.
\]
Strict monotonicity implies $\lambda_p<L$.
\end{proof}


Abbreviate $\lambda=\lambda_p$, and put
\begin{equation}\label{alg:weights-X}
 d_a=\lambda-p_a,\quad w_a=p_a/d_a,\quad
 X_{aa}=1,\quad X_{ab}=2w_aw_b-1\quad(a\ne b).
\end{equation}
Here $X$ is an auxiliary four-point real Gram matrix, distinct from
the Bloch points of an arbitrary competitor.

\begin{lemma}[The auxiliary Gram matrix and its positive kernel]
\label{alg:bloch}
The matrix $X$ is positive semidefinite of rank three. Its kernel is
spanned by the strictly positive vector
\begin{equation}\label{alg:q}
 q_a=\frac{(1-\lambda-p_a)(\lambda-p_a)}
           {2(1-\lambda)(\lambda-2p_a)}.
\end{equation}
It has $\sum q_a=1$, and $q_a<1/2$ for every $a$.
\end{lemma}
\begin{proof}
We have $0<w_a<1$. Set
$D_X=2\operatorname{diag}(1-w_a^2)$, so
$X=D_X+2ww^T-\mathbf1\mathbf1^T$. Write
$A_s=\sum(1-w_a^2)^{-1}$, $B_s=\sum w_a/(1-w_a^2)$,
and $R=\sum(\lambda-2p_a)^{-1}$. Expanding the $2\times2$ determinant
for this rank-two update gives
\[
 \frac{\det X}{\det D_X}=(A_s-3)(1-A_s/2)+B_s^2/2.
\]
The identities
\[
 A_s-B_s=4-1/\lambda,\qquad A_s+B_s=2+\lambda R/2
\]
make $\det X=0$ equivalent to
$(3\lambda-2)(2+\lambda R/2)=8\lambda-5$, which is
$\Phi_p(\lambda)=0$. The matrix $D_X+2ww^T$ is positive definite.
On $\mathbf1^\perp$ its quadratic form remains unchanged after
subtracting $\mathbf1\mathbf1^T$. Thus $X$ has at least three strictly
positive eigenvalues; its remaining eigenvalue is zero.

Every factor in \eqref{alg:q} is positive by Lemma~\ref{alg:root}.
The normalization and kernel identities can be checked without an
eigenvalue computation. For arbitrary $\lambda$ in the displayed
domain, put $\beta=(2\lambda-1)/(2(1-\lambda))$. Division in
$\lambda-2p_a$ gives
\begin{equation}\label{alg:q-residual}
 \sum_aq_a-1=\sum_aw_aq_a-\beta
       =-\frac{\Phi_p(\lambda)}{8(1-\lambda)}.
\end{equation}
Also $2(1-w_a^2)q_a=1-2\beta w_a$. Consequently
$(Xq)_a=(1-2w_a)\Phi_p(\lambda)/(8(1-\lambda))$.
At the root, $\sum q_a=1$ and $Xq=0$.
Realize $X$ by unit vectors $x_a\in\mathbb R^3$. Then $\sum q_ax_a=0$,
so $q_a=\|\sum_{b\ne a}q_bx_b\|\leq1-q_a$.
Equality would force the three other points to coincide at $-x_a$,
contradicting rank three. Hence $q_a<1/2$.
\end{proof}

We use the elementary Bloch representation
$uu^*=(I+x\cdot\sigma)/2$ of a unit spinor $u\in\mathbb C^2$, where
\[
 \sigma_1=\begin{pmatrix}0&1\\1&0\end{pmatrix},\quad
 \sigma_2=\begin{pmatrix}0&-i\\i&0\end{pmatrix},\quad
 \sigma_3=\begin{pmatrix}1&0\\0&-1\end{pmatrix}.
\]
Multiplication and traces of these matrices give
\begin{equation}\label{alg:bloch-identity}
 |u^*v|^2=(1+x\cdot y)/2.
\end{equation}
Conversely a spinor for $x\in S^2$ is
\[
 u(x)=\left(\sqrt{(1+x_3)/2},
       \frac{x_1+i x_2}{\sqrt{2(1+x_3)}}\right)^T
\]
when $x_3\ne-1$, and $(0,1)^T$ at the south pole.

\begin{proposition}[An explicit attaining configuration]\label{alg:construction}
Let $m_a$ be positive integers with total $N$ and admissible
$p_a=m_a/N$. There are $N$ unit vectors in $\mathbb C^{N-2}$ whose Gram
matrix has rank exactly $N-2$ and constant off-diagonal absolute value
\begin{equation}\label{alg:mu-profile}
 \mu_p=\frac1{N\lambda_p-1}.
\end{equation}
To construct them, realize $X$ by real unit vectors $x_a$, choose
corresponding unit spinors $u_a$, and let $W\in\mathbb C^{2\times4}$
have columns $\sqrt{d_a}u_a$. Set
\begin{equation}\label{alg:core}
\begin{aligned}
 B&=W^*W,\qquad A_p=\lambda I_4-B,\\
 D_{\mathrm{mass}}&=\operatorname{diag}(\sqrt{p_a}),\qquad
 C=D_{\mathrm{mass}}^{-1}A_pD_{\mathrm{mass}}^{-1}.
\end{aligned}
\end{equation}
Let $J\in\mathbb R^{N\times4}$ be the isometric block lift, with
$J_{ia}=1/\sqrt{m_a}$ for $i$ in block $a$, and zero otherwise. Then
\begin{equation}\label{alg:signature-Gram}
 S=NJA_pJ^*,\qquad G_p=\frac{N\lambda I_N-S}{N\lambda-1}
\end{equation}
is the required Gram matrix. In particular $S_{ij}=C_{ab}$ for
$i\in a$, $j\in b$.
\end{proposition}
\begin{proof}
The Bloch identity gives $B_{aa}=\lambda-p_a$ and
$|B_{ab}|=\sqrt{d_ad_bw_aw_b}=\sqrt{p_ap_b}$ for $a\ne b$.
The matrix $B$ is positive semidefinite of rank two. Indeed rank one
would force all spinor lines, and hence all Bloch points, to coincide.
Thus $A_p$ has top eigenvalue $\lambda$ of multiplicity exactly two,
while $C_{aa}=1$ and $|C_{ab}|=1$. The spectrum of $S$ is the spectrum
of $NA_p$ together with $N-4$ zeros. Hence $N\lambda I-S$ is positive
semidefinite with exactly two zero eigenvalues. Its diagonal is
$N\lambda-1$, and every off-diagonal entry has modulus one.

An explicit vector factor is obtained as follows.
Choose $R\in\mathbb C^{N\times(N-4)}$ with orthonormal columns
spanning $\operatorname{range}(J)^\perp$. Then
\begin{equation}\label{alg:vector-factor}
 V_p=\frac1{\sqrt{N\lambda-1}}
       \begin{pmatrix}\sqrt{N\lambda}\,R^*\\ \sqrt N\,WJ^*\end{pmatrix}
       \in\mathbb C^{(N-2)\times N}
\end{equation}
satisfies $V_p^*V_p=G_p$, because
$RR^*=I-JJ^*$ and
$N\lambda I-S=N\lambda(I-JJ^*)+NJBJ^*$.
For $N=4$ the first block of rows is empty.
\end{proof}

\begin{remark}[Choices and equivalences]\label{alg:construction-choices}
The construction uses only the specified root, matrix square roots,
and the displayed spinor lift. The first three Bloch vectors can be
obtained by factoring their positive definite principal Gram matrix;
the fourth is then determined by its inner products. Positivity of
that principal matrix follows because the kernel of $X$ has no vector
with a zero coordinate except zero. Any two realizations of $X$ differ
by an orthogonal transformation. A proper rotation corresponds to a
common unitary on the spinors, and an improper rotation to a common
unitary followed by conjugation. One checks the rotation assertion
using $\cos(\theta/2)I-i\sin(\theta/2)n\cdot\sigma$; complex conjugation
reverses the second Bloch coordinate. Independent spinor phases give
diagonal unitary conjugations of $B,C,S,G_p$. Thus for an ordered
profile $G_p$ is unique up to diagonal unitary conjugation and possibly
entrywise conjugation. Permuting the profile permutes the blocks.
\end{remark}

\begin{proposition}[A positive stress with strict phase alignment]
\label{alg:stress}
For the above construction let
\begin{equation}\label{alg:stress-def}
 \ell_a=\frac{2q_a}{d_a}
 =\frac{1-\lambda-p_a}{(1-\lambda)(\lambda-2p_a)},\quad
 L=\operatorname{diag}(\ell_a),\quad T=L-LBL.
\end{equation}
Then $T\succeq0$ has rank two and range $\ker B$.
Choose $U\in\mathbb C^{4\times2}$ with $U^*U=I_2$ and $BU=0$, and set
\begin{equation}\label{alg:Z-K}
 z_a=\frac{U^*e_a}{\sqrt{2p_a}},\qquad
 Z=\sqrt{N/2}\,U^*J^*,\qquad K=\frac{2U^*TU}{\operatorname{tr}T}.
\end{equation}
The matrix $Z$ repeats $z_a$ exactly $m_a$ times and obeys
$ZZ^*=NI_2/2$ and $ZS=N\lambda Z$. Moreover $K\succ0$,
$\operatorname{tr}K=2$, and all entries of $Z^*KZ$ are nonzero with
entrywise phase $S$. Consequently
\begin{equation}\label{alg:candidate-value}
 \Gamma_K(Z):=\sum_{i,j}|z_i^*Kz_j|=N^2\lambda.
\end{equation}
With $K$ held fixed, $Z$ is a critical point of
\begin{equation}\label{alg:Lagrangian}
 \mathcal L(Y)=\Gamma_K(Y)-N\lambda\operatorname{tr}(KYY^*).
\end{equation}
\end{proposition}
\begin{proof}
Since $\sum q_ax_a=0$ and $\sum q_a=1$, the Bloch identity gives
$\sum q_au_au_a^*=I_2/2$, and hence $WLW^*=I_2$.
If $E=WL^{1/2}$, then $EE^*=I_2$, so $E^*E$ is a rank-two orthogonal
projector. Therefore
\[
 T=L^{1/2}(I-E^*E)L^{1/2}\succeq0,\qquad\operatorname{rank}T=2.
\]
Also $WT=WL-(WLW^*)WL=0$. Thus the range of $T$ is $\ker B$,
$U^*TU$ is positive definite, and $\operatorname{tr}(U^*TU)=\operatorname{tr}T$.
The exact entry formulas are
\[
 T_{ab}=\ell_a\ell_b\sqrt{p_ap_b}\,C_{ab}\quad(a\ne b),\qquad
 T_{aa}=\ell_a(1-2q_a)>0.
\]
All factors in these comparisons are strictly positive, by
Lemma~\ref{alg:bloch}. As $UKU^*=2T/\operatorname{tr}T$,
\[
 (Z^*KZ)_{ij}=\frac{T_{ab}}{\operatorname{tr}T\sqrt{p_ap_b}}
       \quad(i\in a,\ j\in b).
\]
This proves the phase assertion even for distinct indices in the same
block. The covariance and eigenvector identities follow from
$J^*J=I$ and $BU=0$. Thus
\[
 \Gamma_K(Z)=\operatorname{tr}(S Z^*KZ)
 =\operatorname{tr}(KZSZ^*)=N^2\lambda.
\]
In a neighborhood of $Z$ all weighted overlaps remain nonzero.
If $C_K(Y)$ is their phase matrix, differentiation with the real
Frobenius inner product gives $\nabla\Gamma_K(Y)=2KYC_K(Y)$.
At $Z$ this is $2KZS=2N\lambda KZ$, exactly the gradient of the
constraint term in \eqref{alg:Lagrangian}.
\end{proof}

\begin{proposition}[Exact balancing]\label{alg:balancing}
Among admissible integer four-weight profiles with total $N$,
$\lambda_p$ is largest and $\mu_p$ is smallest precisely when the four
counts differ by at most one. For every $N\geq10$ the balanced
profile is admissible. If $N\equiv2\pmod4$, then
\begin{equation}\label{alg:residue-two}
 \lambda_{\rm bal}=\frac{1+\sqrt{3-8/N^2}}4,\qquad
 \mu_{\rm bal}=\frac4{N+\sqrt{3N^2-8}-4}.
\end{equation}
\end{proposition}
\begin{proof}
For fixed $\lambda>2/3$, the function
$x\mapsto(\lambda-2x)^{-1}$ is strictly convex. If two counts differ
by at least two, transfer one from the larger to the smaller.
All counts stay positive and below $N/3$, and the sum of reciprocal
denominators strictly decreases. The coefficient $\lambda(3\lambda-2)$
in \eqref{alg:Phi} is positive, so at the old root the new $\Phi$ is
negative. Its unique root is strictly larger. Iteration terminates
at balanced counts, proving also the equality characterization.
For $N\geq10$, $\lceil N/4\rceil<N/3$.
For residue two the weights are $(a,a,b,b)$ with
$a+b=1/2$, $ab=1/16-1/(4N^2)$, and
\[
 \sum_a\frac1{\lambda-2p_a}
       =\frac{4\lambda-2}{\lambda^2-\lambda+4ab}.
\]
The root equation reduces to $2\lambda^2-\lambda-4ab=0$;
its relevant root is $(1+\sqrt{1+32ab})/4$.
\end{proof}


\subsection{The candidate domain needed at the endpoint}\label{app:candidate:domain}
Write $p_a=1/4+\varepsilon_a$, $e=\max|\varepsilon_a|$ and
$\sigma^2=\sum\varepsilon_a^2$. Throughout this subsection $e\le1/120$.
This includes every count profile in Proposition~\ref{prop:localization}
for $N\ge180$. All estimates concern the same alignment of the candidate
and a reference tetrahedron.

The exact scalar identities are
\[
 \Phi_p(\alpha)=4\sum_a\frac{\varepsilon_a^2}{\alpha-1/2-2\varepsilon_a},
 \qquad \Phi'_p(t)=4\left[2+\sum_a\frac{p_a(1-3p_a)}{(t-2p_a)^2}\right].
\]
For $f_t(p)=p(1-3p)/(t-2p)^2$,
\[
 f_t''(p)=\frac{2[t(4-3t)+4p(1-3t)]}{(t-2p)^4}>0
\]
on $2/3\le t\le\alpha<41/60$, $29/120\le p\le31/120$:
the numerator bracket is at least $43/200$. Jensen gives
$\Phi'_p(t)>4568/121>37$, whereas
$4/(\alpha-1/2-2e)<1200/49<74/3$. Therefore
\begin{equation}\label{app:candidate:deficit}
 0\le\alpha-\lambda_p\le\frac23\sigma^2,\qquad
 \frac{109}{100}<b:=3(2\lambda_p-1)<\frac{11}{10},\qquad
 u_0:=\frac34-\lambda_p>\frac1{15}.
\end{equation}
For balanced counts, $N^2\sigma^2=r(4-r)/4$. The odd budget is thus
$1/2$. For residue two the exact formula gives
$N^2(\alpha-\lambda_{\rm bal})=2/(\sqrt3+\sqrt{3-8/N^2})<11/19$
at $N\ge180$. These are the budgets in \eqref{eq:budget}.

\subsection{Exact projector and simultaneous polar estimates}
Let $\Pi=UU^*$ project onto $\ker B$. The two positive eigenvalues of $B$
have sum $4\lambda-1$ and product $D=3\lambda(2\lambda-1)$, whence
\[
 \Pi=I-\frac{(4\lambda-1)B-B^2}{D},\qquad
 \Pi_{aa}=\frac{3\lambda-2+2p_a}{b}.
\]
For complementary pairs $\{a,b\},\{c,d\}$, completion of $U$ to a
unitary matrix and equality of complementary minors give
\[
 \det\Pi[\{a,b\}]=\det(I-\Pi)[\{c,d\}]
 =\frac{\det B[\{c,d\}]}{D}=\frac{\lambda-1+p_a+p_b}{b}.
\]
These identities determine both the lengths $r_a=\|z_a\|$ and the
Bloch Gram matrix $Y_p$ of the normalized candidate spinors, relative
to the tetrahedral Bloch Gram matrix
$Y_0=\frac43I_4-\frac13\mathbf1\mathbf1^T$:
\begin{align}
 r_a^2&=\frac{1+4\varepsilon_a/b}{1+4\varepsilon_a},\label{app:candidate:projector}\\
 (Y_p-Y_0)_{ab}
 &=\frac{-\frac8{3b}(\varepsilon_a+\varepsilon_b)
                +\frac{64}{3b^2}\varepsilon_a\varepsilon_b}
 {(1+4\varepsilon_a/b)(1+4\varepsilon_b/b)}\quad(a\ne b).\label{app:candidate:Ygram}
\end{align}
The diagonal of $Y_p-Y_0$ is zero. The zero-sum identities imply
\[
 G:=\|Y_p-Y_0\|_F\le
 \frac{16/(3b)+112e/(3b^2)}{(1-4e/b)^2}\,\sigma
 \le\frac{11}{2}\sigma\le\frac{11}{120}.
\]
The coefficient at $b=109/100,e=1/120$ is $551200/100489<11/2$.

Here is a polar comparison used also for the actual group means.
For a real $3\times4$ unit-column matrix $X$ and a tetrahedron $V$ with
$V\mathbf1=0$, $VV^T=4I/3$, set $C_0=3XV^T/4$, $c=X\mathbf1/4$.
Then $X=C_0V+c\mathbf1^T$ and, if $G=\|X^TX-V^TV\|_F$,
\begin{equation}\label{app:candidate:polaridentity}
 G^2=\frac{16}{9}\|C_0^TC_0-I\|_F^2
       +\frac{32}{3}\|C_0^Tc\|^2+16\|c\|^4.
\end{equation}
This follows by expanding and using $V\mathbf1=0$. Also
$s_{\min}(C_0)^2\ge1-3G/4$. For the orthogonal polar factor $R$,
\begin{equation}\label{app:candidate:polar}
 \|X-RV\|_F^2\le P_G G^2,\qquad
 P_G=\max\left\{\frac{3/4}{(1+\sqrt{1-3G/4})^2},
                      \frac{3/8}{1-3G/4}\right\}.
\end{equation}
Indeed compare the two nonnegative terms
$\frac43\|C_0-R\|_F^2+4\|c\|^2$ with the first two terms of
\eqref{app:candidate:polaridentity}. At $G\le11/120$,
$s_{\min}(C_0)^2\ge149/160$, so $P_G<41/100$.
An orientation reversal is handled by conjugating the reference spinors.

For unit spinors $v,f$ and their Bloch vectors $x,y$,
\begin{equation}\label{app:candidate:spinor}
 d_{\rm ph}(v,f)^2=2(1-|v^*f|)
 =\frac{\|x-y\|^2}{2(1+|v^*f|)}.
\end{equation}
The total bound $(41/100)(11/120)^2<4[1-(999/1000)^2]$ allows one
common polar alignment and phases with $|v_a^*f_a|>999/1000$.
Moreover
\begin{gather*}
 |r_a^2-1|\le\frac{4e}{11(1-4e)}\le\frac1{319},\qquad
 \frac{499}{500}<r_a<\frac{501}{500},\\
 |r_a-1|\le c_r|\varepsilon_a|,
 \qquad c_r=\frac{4/11}{(1-4e)(1+499/500)}.
\end{gather*}
All these assertions use the larger $1/120$ domain; no smaller-domain
polar floor is extrapolated.

For the actual four count profiles, the following explicit formulas
sharpen the simultaneous errors. Put $\varepsilon_a=\delta_a/N$,
$\sigma=\sqrt C/N$, $b_0=109/100$, $k_0=8/(3b_0)$ and
$q_0=1-4e/b_0$. Write $\Delta_{\mathrm B}=Y_p-Y_0=D_0+F$, where
$(D_0)_{ab}=-8(\varepsilon_a+\varepsilon_b)/(3b)$ off the diagonal.
Exact subtraction in \eqref{app:candidate:Ygram} gives
\[
 F_{ab}=\frac{32}{3b^2}
 \frac{\varepsilon_a^2+4\varepsilon_a\varepsilon_b+\varepsilon_b^2
 +(4/b)(\varepsilon_a+\varepsilon_b)\varepsilon_a\varepsilon_b}
 {(1+4\varepsilon_a/b)(1+4\varepsilon_b/b)}.
\]
Thus $|F_{ab}|\le U_{ab}$, where
\[
 U_{ab}=\frac{32}{3b_0^2N^2q_0^2}
 \left(|\delta_a^2+4\delta_a\delta_b+\delta_b^2|
       +\frac4{b_0N}|(\delta_a+\delta_b)\delta_a\delta_b|\right).
\]
Let $f=(\sum_{a\ne b}U_{ab}^2)^{1/2}/\sigma$.
To see explicitly how this supplies a common polar bound, put
$E=9V\Delta_{\mathrm B}V^T/16$, $z=3V\Delta_{\mathrm B}\mathbf1/16$ and $P_0=(I+E)^{1/2}$.
The aligned Bloch error is
$H=(P_0-I)V+P_0^{-1}z\mathbf1^T$. Its linear parts satisfy
\[
 E_0=\frac{9k}{8}V\diag(\varepsilon_a)V^T,\quad
 z_0=-\frac{3k}{8}V\varepsilon,\quad k=8/(3b),
\]
\[
 \|E_0\|_F=\frac{3\sqrt2 k}{4}\sigma,
 \quad |z_0|=\frac{\sqrt3 k}{4}\sigma.
\]
These follow from $V^TV=\frac43I-\frac13\mathbf1\mathbf1^T$ and
$\sum\varepsilon_a=0$. Set
\[
 K_E=\frac{3\sqrt2}{4}k_0+\frac34f,\quad
 K_z=\frac{\sqrt3}{4}(k_0+f),\quad p_0=\sqrt{1-K_E\sigma},
\]
\[
 R_P=\frac{K_E^2\sigma}{2(1+p_0)^2},\qquad
 R_I=\frac{K_EK_z\sigma}{p_0(1+p_0)}.
\]
The scalar square-root identity and $P_0^{-1}-I$ give these full
nonlinear remainders. The tetrahedral Gram identities, applied separately
to each column, then give
\begin{align*}
 |H_a|/\sigma&\le\frac{\sqrt{3-4\delta_a^2/C}}{b_0}
                          +\frac{13}{16}f+R_P+R_I,\\
 \|H\|_F/\sigma&\le\frac{\sqrt8}{b_0}
                       +\frac{\sqrt{15}}4f+\frac2{\sqrt3}R_P+2R_I.
\end{align*}
For clarity, the linear point norm equals
$\sqrt{3\sigma^2-4\varepsilon_a^2}/b$; its aggregate norm is
$\sqrt8\,\sigma/b$. The nonlinear $F$ contribution is bounded by
$13\|F\|_F/16$ per column and $\sqrt{15}\|F\|_F/4$ in aggregate,
using $E_F=9VFV^T/16$, $z_F=3VF\mathbf1/16$ and
$\|E_F\|_F\le3\|F\|_F/4$, $|z_F|\le\sqrt3\|F\|_F/4$.
The two arrays $E_FV/2$ and $z_F\mathbf1^T$ are orthogonal, giving
the aggregate squared coefficient $15/16$; the point coefficient is
$3/8+\sqrt3/4<13/16$. This explains the paired, rather than independently
maximized, column dependence.

With $\phi=(501/500)/[2(1+999/1000)]$ and
$R_*=(13/16)f+R_P+R_I$, define
\begin{align*}
 C_{\mathrm{pt},a}&=c_r^2\delta_a^2/C+
      \phi[\sqrt{3-4\delta_a^2/C}/b_0+R_*]^2,\\
 C_{\mathrm{pt}}&=\max_a C_{\mathrm{pt},a},\\
 B_p&=(1/4+e)\left\{c_r^2+
 \phi[\sqrt8/b_0+(\sqrt{15}/4)f+(2/\sqrt3)R_P+2R_I]^2\right\}.
\end{align*}
In this same alignment they prove
\begin{equation}\label{app:candidate:PB}
 \max_a\|z_a-f_a\|^2\le C_{\mathrm{pt}}C/N^2,
 \qquad \sum_a p_a\|z_a-f_a\|^2\le B_p C/N^2.
\end{equation}
All divisions are guarded by $p_0>0$. On the full domain the generic
bound $f<61/100$ implies $K_E<33/8$, $K_z<77/32$, hence positivity;
actual profile bounds are sharper. Both signs and every permutation
are included in the absolute polynomials $U_{ab}$.

Finally the stress norm has a basis-free estimate. Let $j_a=1/\ell_a$
and $H=U^*\diag(j_a)U$. The identity $TL^{-1}\Pi=\Pi$ gives
$U^*TU=H^{-1}$ and $\|K-I\|=\operatorname{gap}(H)/\tr H$.
Writing $j_c=j(1/4,\lambda)$, exact division gives
\[
 j_a/j_c=1+C_0\varepsilon_a+
             C_0\varepsilon_a^2/(u_0-\varepsilon_a),\quad
 C_0=\frac{3\lambda-2}{(\lambda-1/2)u_0},\quad0\le C_0\le4.
\]
For $E=U^*\diag(\varepsilon_a)U$, the projector minors give
\[
 \tr E=2\sigma^2/b,\qquad
 \operatorname{gap}(E)^2=\sigma^2/3+
             4\sum\varepsilon_a^3/b+4\sigma^4/b^2.
\]
The remaining compression is positive semidefinite and has spectral
gap at most $4e^2/(u_0-e)$, whereas $\tr(H/j_c)\ge2$. Consequently
\begin{equation}\label{app:candidate:kappa}
 \kappa:=\|K-I\|\le
 2\sigma\sqrt{1/3+4e/b_0+4\sigma^2/b_0^2}
                         +\frac{2e^2}{1/15-e}<\frac{11}4e.
\end{equation}
The last endpoint check is
$4\sqrt{39013/106929}+2/7<11/4$. The radial and spinor estimates also give a generic bound $\eta$
for the candidate point error $\max_a\|z_a-f_a\|$:
$\eta^2\le[500/(319\cdot999)]^2+
\phi(41/100)(11/120)^2<1/1024$.
The profile-dependent choice of $\eta$ is specified in
\eqref{app:local:polar-input}.
Thus the candidate point error and $\kappa$ are below $1/32$.
For a fixed count profile all
analytic upper coefficients above improve as $N$ increases. Freezing
valid endpoint enclosures therefore proves the whole subsequent range;
no claim about monotonicity of rounded program outputs is needed.
